\section{Prospective Benchmark Results}
\label{sec:prospective}

\subsection{Scope and Audit Trail}

The final benchmark contains 11 datasets. Cora, CiteSeer, PubMed, and Coauthor-CS represent citation and co-authorship graphs; Chameleon, Squirrel, Actor, Wisconsin, Cornell, Roman-empire, and Amazon-ratings broaden the structural regimes. Each dataset uses ten fixed seeds, with 60\% of nodes for training, 20\% for validation, and 20\% for testing. The run yields 770 selected-model records and 110 diagnostic records: every unit contains the MLP and six graph architectures, four tuning trials per architecture, explicit split and environment provenance, and one test evaluation after validation selection. No final-run unit fails.

The execution history contains two documented amendments. First, the v1 feasibility run stops before model training on Texas because one class contains a single node and cannot populate the frozen three-way split. Before comparative assembly or evaluation, an outcome-independent eligibility amendment excludes Texas, assigns a new run identifier, and restarts all remaining datasets from zero without changing models, splits, seeds, thresholds, metrics, or inference. Second, after a blinded wall-clock timeout on Squirrel, an execution-only amendment increases the per-dataset limit from 30 to 90 minutes. No comparative outcome is inspected, and all scientific settings remain unchanged. The completed v2 artifacts retain the amendments and exclude the aborted v1 records from analysis.

\subsection{Primary Policy Outcomes}

Graph is the target action in 89 of 110 units, so a selector must be compared with the trivial always-graph policy rather than judged by raw accuracy alone. Table~\ref{tab:prospective_policy_results} reports every frozen policy on identical units. The historical combined diagnostic reaches 55.5\% selection accuracy at 68.2\% coverage. Its covered decisions incur 2.64 percentage points of regret, but the common MLP fallback on 35 abstentions raises full-set regret to 7.46 points.

\begin{table}[t]
\centering
\caption{Prospective Graph-vs-MLP Policy Results on 110 Dataset--Seed Units. Regret is reported in percentage points; selection accuracy uses the common MLP fallback.}
\label{tab:prospective_policy_results}
\small
\begin{tabular}{@{}lccc@{}}
\toprule
Policy & Accuracy (\%) & Coverage (\%) & Regret (pp) \\
\midrule
Combined & 55.5 & 68.2 & 7.46 \\
Always-graph & 80.9 & 100.0 & 0.26 \\
Always-MLP & 19.1 & 100.0 & 9.69 \\
Random 50/50 & 50.0 & 100.0 & 4.97 \\
Homophily-only & 55.5 & 100.0 & 7.46 \\
Degree-only & 37.3 & 100.0 & 6.24 \\
Homophily+degree & 46.4 & 100.0 & 5.80 \\
Two-hop-only & 28.2 & 100.0 & 8.27 \\
Validation selection & 91.8 & 100.0 & 0.22 \\
\bottomrule
\end{tabular}
\end{table}

Always-graph attains 80.9\% selection accuracy and 0.26 points of regret. Direct validation selection attains 91.8\% and 0.22 points, respectively. Homophily-only has the same full-set accuracy and regret as the combined policy after fallback, while degree-only, homophily-plus-degree, and two-hop-only also fail to improve on always-graph. The result is therefore not that a different fixed threshold rescues the original idea; under this portfolio and target, the evaluated low-dimensional rules provide no stable incremental decision value.

\subsection{Paired Inference and Interpretation}

Relative to the combined diagnostic, the dataset-level mean regret difference is $-7.20$ points for always-graph. Its 95\% bootstrap interval is $[-13.17,-1.70]$, with a raw sign-flip $p$-value of $0.015625$. The difference for validation selection is $-7.24$ points, with interval $[-13.18,-1.78]$ and raw $p=0.03125$. After Holm correction across all eight comparisons, the adjusted values are $0.125$ and $0.21875$. We therefore do not claim confirmatory superiority for either reference policy, and we do not treat the non-significant adjusted tests as equivalence.

The operational conclusion is narrower and does not depend on a significance claim in the opposite direction. The combined rule fails its predeclared incremental-value criterion because it does not reduce regret relative to simple policies and its uncertainty does not establish an advantage. Selection accuracy alone would hide this result: the target imbalance makes always-graph strong, while regret shows that the diagnostic's errors are also substantially more costly. Section~\ref{sec:edge_intervention} next tests whether this failure can be explained by graph structure being uninformative.
